\section{子任务 3：第 2 章 Context-Free Languages}

\subsection{核心知识点}

\begin{itemize}[leftmargin=2em]
  \item 上下文无关文法 CFG：用递归生成方式描述语言。
  \item 歧义：同一个串可能存在多棵语法树，说明“结构解释”不唯一。
  \item Chomsky normal form：把文法规整化，便于算法与证明。
  \item 下推自动机 PDA：有限控制加一个栈，专门处理嵌套结构。
  \item CFG 与 PDA 等价：生成与识别是同一个结构现象的两种侧面。
  \item 上下文无关语言的 pumping lemma：利用语法树中变量重复得到可泵部分。
  \item 确定型上下文无关语言、LR(k) 文法与分析：讨论哪些嵌套结构可以被确定地解析。
\end{itemize}

\subsection{典型问题与证明套路}

\begin{enumerate}[leftmargin=2em]
  \item \textbf{证明一个语言是 CFL}：直接写 CFG，核心是找到语言的递归拆分方式。
  \item \textbf{如果语言带有括号匹配、成对嵌套、递归块结构}：优先设计 PDA，用栈保存“尚未闭合的义务”。
  \item \textbf{证明 CFG 与 PDA 等价}：从语法树导出栈操作，或让 PDA 模拟最左推导。
  \item \textbf{证明某文法算法可行}：先化成 CNF，让推导只剩少数标准形态，再做动态规划或归纳。
  \item \textbf{证明不是 CFL}：使用 pumping lemma，选取要求多重同步约束的长串，例如多段长度必须同时相等的串，然后逐类分析被泵部分落在哪一段。
  \item \textbf{证明一个语言需要确定性还是非确定性}：分析解析时是否需要“猜”，以及有限前瞻能否消除这种猜测。
\end{enumerate}

\subsection{证明背后的哲学}

这一章的哲学是：\textbf{一个栈足以表达递归和嵌套，但通常不足以表达多个彼此独立的无界计数。}
因此，像括号匹配、函数调用、块结构这些“后进先出”的关系，非常自然地落在 CFL 内；而要求三段同时同步、两个方向同时精确记账的问题，常常超出 PDA 的能力。

另一个重要哲学是：\textbf{语法不是装饰，而是一种计算。} 语法树的形状本身就编码了识别过程中的关键中间信息。

\subsection{本章精华}

\begin{itemize}[leftmargin=2em]
  \item 正则语言解决“局部模式”，上下文无关语言解决“嵌套模式”。
  \item 从 CFG 到 PDA 再到 LR(k)，展示了从纯理论到编译原理的自然桥梁。
\end{itemize}

\newpage
